%==============================================================================
%  UG PROJECT REPORT  --  Chandigarh University format
%  Compile with: pdfLaTeX  (Overleaf: Menu -> Compiler -> pdfLaTeX)
%  Compile 2-3 times so the Table of Contents / List of Figures fill in.
%
%  IMAGES TO UPLOAD (drag into the Overleaf project, same folder as this file):
%     cu-logo.png     <- Chandigarh University logo (used on cover + certificate)
%  If the file is missing the document still compiles and shows a grey box.
%==============================================================================
\documentclass[12pt,a4paper]{report}

%------------------------- page geometry (1 inch all round) -------------------
\usepackage[a4paper,top=1in,bottom=1in,left=1in,right=1in]{geometry}
% For a binding margin instead, use:
% \usepackage[a4paper,top=1in,bottom=1in,left=1.25in,right=1in]{geometry}

%------------------------- fonts: Times New Roman clone -----------------------
\usepackage[T1]{fontenc}
\usepackage[utf8]{inputenc}
\usepackage{mathptmx}          % Nimbus Roman = metrically identical to Times New Roman

%------------------------- spacing, layout, tools -----------------------------
\usepackage{setspace}
\usepackage{graphicx}
\usepackage{xcolor}
\usepackage{titlesec}
\usepackage{tocloft}
\usepackage[font={bf,footnotesize},labelsep=colon,justification=centering,skip=6pt]{caption}
\usepackage{enumitem}
\usepackage{array}
\usepackage{longtable}
\usepackage{float}
\usepackage{multirow}
\usepackage{amsmath,amssymb}
\usepackage[hidelinks,bookmarksnumbered]{hyperref}   % keep hyperref last

\onehalfspacing                        % 1.5 line spacing for the general text
\setlength{\parindent}{0pt}            % Word-style: no first-line indent
\setlength{\parskip}{6pt}
\setlist{nosep,topsep=4pt,itemsep=3pt,leftmargin=1.2cm}
\renewcommand{\arraystretch}{1.3}
% ragged-right paragraph column for tables:  L{3cm}  instead of  p{3cm}
\newcolumntype{L}[1]{>{\raggedright\arraybackslash}p{#1}}

% Word does not hyphenate by default -- switch it off so the justification matches
\frenchspacing          % Word-style: no extra space after a full stop
\hyphenpenalty=10000
\exhyphenpenalty=10000
\tolerance=9999
\emergencystretch=3em

%------------------------- exact font sizes from the format sheet -------------
\newcommand{\sixteen}{\fontsize{16}{19.2}\selectfont}   % chapter names
\newcommand{\fourteen}{\fontsize{14}{21}\selectfont}    % main headings
\newcommand{\twelve}{\fontsize{12}{18}\selectfont}      % body / sub-headings
\newcommand{\eighteen}{\fontsize{18}{27}\selectfont}    % report title

%------------------------- CHAPTER: 16pt, bold, centered ----------------------
\titleformat{\chapter}[display]
  {\filcenter\bfseries\sixteen}
  {CHAPTER \thechapter.}
  {6pt}
  {}
\titlespacing*{\chapter}{0pt}{-24pt}{22pt}

% unnumbered chapters (REFERENCES, APPENDIX, USER MANUAL ...)
\titleformat{name=\chapter,numberless}[display]
  {\filcenter\bfseries\sixteen}{}{0pt}{}
\titlespacing*{name=\chapter,numberless}{0pt}{-24pt}{22pt}

%------------------------- HEADING: 14pt, bold, left ---------------------------
\titleformat{\section}[hang]
  {\bfseries\fourteen}
  {\makebox[3.4em][l]{\thesection.}}{0pt}{}
\titlespacing*{\section}{0pt}{16pt}{8pt}

%------------------------- SUB-HEADING: 12pt, bold, left -----------------------
\titleformat{\subsection}[hang]
  {\bfseries\twelve}
  {\makebox[3.4em][l]{\thesubsection.}}{0pt}{}
\titlespacing*{\subsection}{0pt}{12pt}{6pt}

\titleformat{\subsubsection}[hang]
  {\bfseries\twelve}
  {\makebox[3.9em][l]{\thesubsubsection.}}{0pt}{}
\titlespacing*{\subsubsection}{0pt}{10pt}{6pt}

\setcounter{secnumdepth}{3}
\setcounter{tocdepth}{2}

%------------------------- TABLE OF CONTENTS styling ---------------------------
\renewcommand{\contentsname}{TABLE OF CONTENTS}
\renewcommand{\cfttoctitlefont}{\bfseries\sixteen}
\setlength{\cftbeforetoctitleskip}{0pt}
\setlength{\cftaftertoctitleskip}{22pt}

% "Chapter 1. ....... 4"
\renewcommand{\cftchappresnum}{Chapter\ }
\renewcommand{\cftchapaftersnum}{.}
\setlength{\cftchapnumwidth}{6.2em}
\renewcommand{\cftchapfont}{\bfseries}
\renewcommand{\cftchappagefont}{\bfseries}
\renewcommand{\cftchapleader}{\bfseries\cftdotfill{\cftdotsep}}
\renewcommand{\cftchapdotsep}{\cftdotsep}
\setlength{\cftbeforechapskip}{8pt}

\renewcommand{\cftsecfont}{\bfseries}
\renewcommand{\cftsecpagefont}{\bfseries}
\renewcommand{\cftsecleader}{\bfseries\cftdotfill{\cftdotsep}}
\setlength{\cftsecnumwidth}{3.2em}
\setlength{\cftbeforesecskip}{6pt}

\renewcommand{\cftsubsecfont}{\bfseries}
\renewcommand{\cftsubsecleader}{\bfseries\cftdotfill{\cftdotsep}}
\renewcommand{\cftsubsecpagefont}{\bfseries}
\setlength{\cftsubsecnumwidth}{4.0em}
\setlength{\cftbeforesubsecskip}{6pt}

%------------------------- LIST OF FIGURES / TABLES styling --------------------
\renewcommand{\listfigurename}{List of Figures}
\renewcommand{\cftloftitlefont}{\bfseries\sixteen}
\setlength{\cftbeforeloftitleskip}{0pt}
\setlength{\cftafterloftitleskip}{22pt}
\renewcommand{\cftfigpresnum}{Figure\ }
\renewcommand{\cftfigaftersnum}{:}
\setlength{\cftfignumwidth}{5.6em}
\renewcommand{\cftfigfont}{\bfseries}
\renewcommand{\cftfigpagefont}{\bfseries}
\renewcommand{\cftfigleader}{\bfseries\cftdotfill{\cftdotsep}}
\setlength{\cftbeforefigskip}{8pt}

\renewcommand{\listtablename}{List of Tables}
\renewcommand{\cftlottitlefont}{\bfseries\sixteen}
\setlength{\cftbeforelottitleskip}{0pt}
\setlength{\cftafterlottitleskip}{22pt}
\renewcommand{\cfttabpresnum}{Table\ }
\renewcommand{\cfttabaftersnum}{:}
\setlength{\cfttabnumwidth}{5.2em}
\renewcommand{\cfttabfont}{\bfseries}
\renewcommand{\cfttabpagefont}{\bfseries}
\renewcommand{\cfttableader}{\bfseries\cftdotfill{\cftdotsep}}
\setlength{\cftbeforetabskip}{8pt}

% centre the three list titles (tocloft's own \hfill trick loses to \parfillskip)
\makeatletter
\renewcommand{\@cftmaketoctitle}{%
  \addpenalty\@secpenalty \vspace*{\cftbeforetoctitleskip}\@cftpagestyle
  {\interlinepenalty\@M {\centering\cfttoctitlefont\contentsname\par}%
   \cftmarktoc \nobreak \vskip\cftaftertoctitleskip \@afterheading}}
\renewcommand{\@cftmakeloftitle}{%
  \addpenalty\@secpenalty \vspace*{\cftbeforeloftitleskip}\@cftpagestyle
  {\interlinepenalty\@M {\centering\cftloftitlefont\listfigurename\par}%
   \cftmarklof \nobreak \vskip\cftafterloftitleskip \@afterheading}}
\renewcommand{\@cftmakelottitle}{%
  \addpenalty\@secpenalty \vspace*{\cftbeforelottitleskip}\@cftpagestyle
  {\interlinepenalty\@M {\centering\cftlottitlefont\listtablename\par}%
   \cftmarklot \nobreak \vskip\cftafterlottitleskip \@afterheading}}
\makeatother

%------------------------- helper: logo with safe fallback ---------------------
\newcommand{\culogo}[1][3.15in]{%
  \IfFileExists{cu-logo.png}
    {\includegraphics[width=#1]{cu-logo.png}}
    {\fbox{\parbox[c][1.1in][c]{#1}{\centering\small\itshape
      upload \texttt{cu-logo.png}\\to the project folder}}}%
}

%------------------------- helper: placeholder figure box ----------------------
% Replace \phbox{...} with \includegraphics[width=...]{yourfile.png}
\newcommand{\phbox}[2][0.72\textwidth]{%
  \fbox{\parbox[c][#2][c]{#1}{\centering\fontsize{34}{40}\selectfont Figure}}%
}

%------------------------- front-matter page helper ----------------------------
\newcommand{\frontchapter}[1]{%
  \chapter*{#1}\phantomsection\addcontentsline{toc}{chapter}{#1}%
}

%==============================================================================
%  >>>>>>>>>>>>>>>>>>>>>  EDIT YOUR DETAILS HERE  <<<<<<<<<<<<<<<<<<<<<<<<<<<<
%==============================================================================
\newcommand{\ProjectTitle}{IOT-BASED SMART ENERGY MONITORING SYSTEM}
\newcommand{\CandidateOne}{John Doe (21BEE1001)}
\newcommand{\CandidateTwo}{Jane Doe (21BEE1002)}
\newcommand{\DegreeName}{BACHELOR OF ENGINEERING}
\newcommand{\BranchName}{ELECTRONICS ENGINEERING}
\newcommand{\MonthYear}{MAY 2026}
\newcommand{\HoDName}{Dr. A. B. Sharma}
\newcommand{\SupervisorName}{Dr. C. D. Verma}
\newcommand{\SupervisorDesignation}{Associate Professor}
\newcommand{\DeptName}{Department of Electronics \& Communication Engineering}
%==============================================================================

\begin{document}

%==============================================================================
%  COVER / TITLE PAGE
%==============================================================================
\pagestyle{empty}
\pagenumbering{roman}

\begingroup
\setlength{\parskip}{0pt}
\singlespacing
\begin{center}
\vspace*{0.20in}

{\bfseries\eighteen \ProjectTitle\par}

\vspace{0.50in}
{\bfseries\fourteen A PROJECT REPORT\par}

\vspace{0.50in}
{\bfseries\itshape\fourteen Submitted by\par}

\vspace{0.45in}
{\bfseries\sixteen \CandidateOne\par}
\vspace{0.12in}
{\bfseries\sixteen \CandidateTwo\par}

\vspace{0.50in}
{\bfseries\itshape\fourteen in partial fulfillment for the award of the degree of\par}

\vspace{0.50in}
{\bfseries\sixteen \DegreeName\par}

\vspace{0.25in}
{\bfseries\fourteen IN\par}

\vspace{0.25in}
{\fourteen \BranchName\par}

\vspace{0.45in}
\culogo

\vspace{0.12in}
{\bfseries\fourteen Chandigarh University\par}

\vspace{0.35in}
{\fourteen \MonthYear\par}

\end{center}
\endgroup
\clearpage

%==============================================================================
%  BONAFIDE CERTIFICATE
%==============================================================================
\begin{center}
\vspace*{0.1in}
\culogo

\vspace{0.35in}
{\bfseries\sixteen BONAFIDE CERTIFICATE\par}
\end{center}

\vspace{0.40in}

{\fourteen
Certified that this project report \textbf{``\ProjectTitle''} is the bonafide
work of \textbf{``\CandidateOne{} and \CandidateTwo''} who carried out the
project work under my/our supervision.\par}

\vspace{0.55in}

\begin{singlespace}
\noindent
\begin{minipage}[t]{0.47\textwidth}
\raggedright
\fourteen
$<<$Signature of the HoD$>>$\\[0.45in]
\textbf{SIGNATURE}\\[0.18in]
\HoDName\\[0.10in]
\textbf{HEAD OF THE DEPARTMENT}\\[0.10in]
\DeptName
\end{minipage}
\hfill
\begin{minipage}[t]{0.47\textwidth}
\raggedright
\fourteen
$<<$Signature of the Supervisor$>>$\\[0.45in]
\textbf{SIGNATURE}\\[0.18in]
\SupervisorName\\[0.10in]
\textbf{SUPERVISOR}\\[0.10in]
\SupervisorDesignation\\[0.10in]
\DeptName
\end{minipage}
\end{singlespace}

\vspace{0.50in}

{\fourteen Submitted for the project viva-voce examination held on
\rule{1.6in}{0.4pt}\par}

\vspace{0.55in}

\begin{singlespace}
\noindent
\begin{minipage}[t]{0.47\textwidth}
\raggedright\bfseries\fourteen INTERNAL EXAMINER
\end{minipage}
\hfill
\begin{minipage}[t]{0.47\textwidth}
\raggedright\bfseries\fourteen EXTERNAL EXAMINER
\end{minipage}
\end{singlespace}

\clearpage

%==============================================================================
%  ACKNOWLEDGEMENT
%==============================================================================
\begin{center}
\vspace*{0.1in}
{\bfseries\sixteen ACKNOWLEDGEMENT\par}
\end{center}
\vspace{0.3in}

We would like to express our sincere gratitude to our supervisor,
\SupervisorName, \SupervisorDesignation, \DeptName, Chandigarh University,
for the continuous guidance, patience and encouragement extended throughout
the course of this project work.

We are equally thankful to \HoDName, Head of the Department, for providing the
laboratory facilities and the academic environment that made this work
possible. Our thanks also go to the faculty members and technical staff of the
department for their timely help during fabrication and testing.

Finally, we thank our families and friends for their constant support and
motivation.

\vspace{0.6in}
\hfill\begin{minipage}{0.42\textwidth}
\raggedright
\CandidateOne\\
\CandidateTwo
\end{minipage}

\clearpage

%==============================================================================
%  TABLE OF CONTENTS
%==============================================================================
\tableofcontents
\addtocontents{toc}{\protect\thispagestyle{empty}}
\clearpage

%==============================================================================
%  From here the roman page numbers become visible: i, ii, iii ...
%==============================================================================
\pagestyle{plain}
\setcounter{page}{1}

\phantomsection\addcontentsline{toc}{chapter}{List of Figures}
\listoffigures
\clearpage

\phantomsection\addcontentsline{toc}{chapter}{List of Tables}
\listoftables
\clearpage

%==============================================================================
%  ABSTRACT
%==============================================================================
\frontchapter{ABSTRACT}

Rising electricity demand in residential and small commercial buildings has made
appliance-level energy visibility a practical necessity rather than a
convenience. Conventional energy meters report only aggregate consumption at the
end of a billing cycle, which gives the consumer no actionable insight into
where the energy is actually going.

This project presents the design and implementation of an IoT-based smart energy
monitoring system capable of measuring voltage, current, power factor and
cumulative energy at the individual load level. The system uses a
non-invasive current transformer for sensing, an ESP32 microcontroller for
acquisition and edge processing, and MQTT over Wi-Fi for telemetry to a
cloud dashboard. A moving-window algorithm is applied at the edge to reduce
transmitted payload by grouping stable readings.

The prototype was tested against a calibrated reference meter across a load
range of 40~W to 2200~W. Measured energy readings tracked the reference within
$\pm$2.1\% and the dashboard latency remained under 3~seconds. The results
confirm that a low-cost, retrofittable monitoring node can deliver
metering-grade visibility without modification of existing wiring.

\vspace{12pt}
\textbf{Keywords:} Internet of Things, energy monitoring, ESP32, MQTT,
non-invasive sensing, smart metering.

\clearpage

%==============================================================================
%  GRAPHICAL ABSTRACT
%==============================================================================
\frontchapter{GRAPHICAL ABSTRACT}

\vspace{0.4in}
\begin{center}
\phbox[0.85\textwidth]{3.2in}
\end{center}

\clearpage

%==============================================================================
%  ABBREVIATIONS
%==============================================================================
\frontchapter{ABBREVIATIONS}

\begin{singlespace}
\begin{center}
\begin{tabular}{@{}p{3.5cm}p{9.5cm}@{}}
\textbf{ADC}   & Analog to Digital Converter \\
\textbf{API}   & Application Programming Interface \\
\textbf{CT}    & Current Transformer \\
\textbf{IoT}   & Internet of Things \\
\textbf{MQTT}  & Message Queuing Telemetry Transport \\
\textbf{PCB}   & Printed Circuit Board \\
\textbf{RMS}   & Root Mean Square \\
\textbf{SoC}   & System on Chip \\
\textbf{THD}   & Total Harmonic Distortion \\
\end{tabular}
\end{center}
\end{singlespace}

\clearpage

%==============================================================================
%  SYMBOLS
%==============================================================================
\frontchapter{SYMBOLS}

\begin{singlespace}
\begin{center}
\begin{tabular}{@{}p{3.5cm}p{9.5cm}@{}}
$V_{rms}$   & Root mean square voltage (V) \\
$I_{rms}$   & Root mean square current (A) \\
$P$         & Active power (W) \\
$S$         & Apparent power (VA) \\
$\cos\phi$  & Power factor \\
$E$         & Cumulative energy (kWh) \\
$f$         & Supply frequency (Hz) \\
$\varepsilon$ & Relative measurement error (\%) \\
\end{tabular}
\end{center}
\end{singlespace}

\clearpage

%==============================================================================
%  MAIN TEXT
%==============================================================================
\pagenumbering{arabic}
\setcounter{page}{1}

%------------------------------------------------------------------------------
\chapter{INTRODUCTION}

\section{Client Identification/Need Identification/Identification of relevant Contemporary issue}

\begin{itemize}
  \item Justify that the issue at hand exists through statistics and documentation.
  \item It is a problem that someone needs resolution for (consultancy problem).
  \item The need is justified through a survey or reported after a survey.
  \item Relevant contemporary issue documented in reports of some agencies.
\end{itemize}

According to the Central Electricity Authority, domestic consumption accounted
for roughly one quarter of national electricity sales in the last reported
year, and per-capita consumption has grown steadily over the past decade. A
campus survey of 120 households conducted for this project found that 78\% of
respondents could not name their single largest domestic load, and 64\%
reported that their monthly bill regularly surprised them.

The client for this work is the campus facilities office, which maintains
several hostel blocks with no sub-metering. Their stated need is per-block and
per-appliance visibility that can be retrofitted without rewiring.

\section{Identification of Problem}

Occupants and facility managers currently have no mechanism to attribute energy
consumption to individual loads within a building, and therefore cannot
identify, prioritise or verify any reduction measure they attempt.

\section{Identification of Tasks}

The work was decomposed into the following tasks:

\begin{enumerate}
  \item Survey of existing metering practice and literature on non-invasive sensing.
  \item Selection of sensing elements and specification of accuracy targets.
  \item Design of the acquisition hardware and the edge firmware.
  \item Design of the communication stack and cloud dashboard.
  \item Calibration against a reference meter and error characterisation.
  \item Validation across the full rated load range and documentation.
\end{enumerate}

\section{Timeline}

The project was executed over two semesters. The schedule is represented as a
Gantt chart in Figure~\ref{fig:gantt}, showing overlap between the firmware and
dashboard workstreams during weeks 9 to 14.

\begin{figure}[H]
  \centering
  \phbox{2.6in}
  \caption{Gantt chart of the project schedule}
  \label{fig:gantt}
\end{figure}

\section{Organization of the Report}

Chapter~2 reviews the reported literature and closes with the formal problem
definition and objectives. Chapter~3 covers concept generation, the design
constraints considered, the two alternative design flows evaluated and the
justification for the selected design. Chapter~4 presents the implementation
using modern engineering tools together with the test and validation results.
Chapter~5 concludes the work and identifies the way ahead.

%------------------------------------------------------------------------------
\chapter{LITERATURE REVIEW/BACKGROUND STUDY}

\section{Timeline of the reported problem}

Interest in appliance-level energy disaggregation dates to the late 1980s, when
non-intrusive load monitoring was first proposed as an alternative to
per-circuit metering. Commercial interest accelerated after 2010 with the
availability of low-cost wireless microcontrollers, and standardisation efforts
around lightweight messaging protocols followed shortly after.

\section{Proposed solutions}

Reported approaches fall into three broad families: single-point disaggregation
from the service entry, distributed plug-level metering, and hybrid schemes
that combine a coarse service-entry meter with a small number of distributed
nodes.

\subsection{Single-point disaggregation}

These systems infer appliance states from transient signatures observed at the
service entry. They require only one sensing point but depend heavily on
training data and degrade when several similar loads operate simultaneously.

\subsection{Distributed plug-level metering}

Each outlet carries its own meter. Accuracy is high and no inference is
required, but cost scales linearly with the number of monitored points and
installation is intrusive.

\section{Bibliometric analysis}

A comparison of the reviewed approaches on key features, effectiveness and
drawbacks is summarised in Table~\ref{tab:biblio}.

\begin{table}[H]
  \centering
  \caption{Comparison of reported monitoring approaches}
  \label{tab:biblio}
  \begin{tabular}{|L{3.2cm}|L{3.4cm}|L{3.0cm}|L{3.0cm}|}
    \hline
    \textbf{Approach} & \textbf{Key feature} & \textbf{Effectiveness} & \textbf{Drawback} \\
    \hline
    Single-point NILM & One sensing point & 70--85\% accuracy & Needs training data \\
    \hline
    Plug-level meters & Direct measurement & $>$98\% accuracy & High unit cost \\
    \hline
    Hybrid scheme & Few nodes + inference & 90--95\% accuracy & Moderate complexity \\
    \hline
    Smart main meter & Utility grade & Billing accurate & No appliance detail \\
    \hline
  \end{tabular}
\end{table}

\section{Review Summary}

The literature indicates that a small number of well-placed distributed nodes
recovers most of the accuracy of full plug-level metering at a fraction of the
cost. This project therefore adopts the hybrid position, with non-invasive
sensing to keep installation retrofittable.

\section{Problem Definition}

To design, build and validate a retrofittable monitoring node that reports
voltage, current, active power and cumulative energy for an individual load to
a cloud dashboard, with measurement error within $\pm$3\% over a 40~W to
2200~W range. The work does not include appliance classification, utility
billing integration, or three-phase operation.

\section{Goals/Objectives}

\begin{itemize}
  \item Achieve energy measurement error within $\pm$3\% against a reference meter.
  \item Keep end-to-end dashboard latency below 5~seconds.
  \item Restrict per-node bill of materials to under Rs.~1,500.
  \item Install without breaking or rewiring the existing circuit.
  \item Sustain 72~hours of continuous unattended logging without data loss.
\end{itemize}

%------------------------------------------------------------------------------
\chapter{DESIGN FLOW/PROCESS}

\section{Evaluation \& Selection of Specifications/Features}

The features identified in the literature were critically evaluated and scored
against the client need. Retained features were non-invasive current sensing,
true RMS computation, local buffering during network loss, and over-the-air
firmware update. Rejected features were appliance classification and on-device
display, both of which added cost without addressing the stated need.

\section{Design Constraints}

\begin{itemize}
  \item \textbf{Regulations:} enclosure and clearances per IS~732 wiring practice.
  \item \textbf{Economic:} node cost capped by the client at Rs.~1,500.
  \item \textbf{Environmental:} lead-free assembly; no primary cells.
  \item \textbf{Health and Safety:} no exposed mains; reinforced isolation.
  \item \textbf{Manufacturability:} single-sided PCB, all through-hole parts.
  \item \textbf{Ethical and Social:} consumption data anonymised at source.
\end{itemize}

\section{Analysis and Feature finalization subject to constraints}

The isolation requirement removed the direct-shunt current sensing option. The
cost cap removed the dedicated metering IC in favour of computing RMS values in
firmware. The manufacturability constraint fixed the enclosure size, which in
turn set the maximum current transformer aperture.

\section{Design Flow}

Two alternative designs were carried forward for comparison.

\subsection{Design~I: Metering IC based node}

A dedicated energy metering IC performs all measurement and reports over SPI to
a low-cost microcontroller. Accuracy is excellent and firmware effort is small,
but the IC alone consumes most of the cost budget.

\subsection{Design~II: Microcontroller with firmware RMS}

The ESP32 samples the sensor chain directly through its internal ADC and
computes RMS quantities in firmware. Cost is low and the design is flexible,
at the price of a calibration routine and higher firmware complexity. The
block diagram of this alternative is shown in Figure~\ref{fig:block}.

\begin{figure}[H]
  \centering
  \phbox{2.9in}
  \caption{Block diagram of the proposed monitoring node}
  \label{fig:block}
\end{figure}

\section{Design selection}

The two alternatives were scored against the constraint set, as shown in
Table~\ref{tab:designsel}. Design~II was selected because it meets the accuracy
target after calibration while remaining within the cost cap, and because it
leaves headroom for the edge-side payload reduction algorithm.

\begin{table}[H]
  \centering
  \caption{Comparison and selection of the design alternatives}
  \label{tab:designsel}
  \begin{tabular}{|L{4.2cm}|L{3.6cm}|L{3.6cm}|}
    \hline
    \textbf{Criterion} & \textbf{Design~I} & \textbf{Design~II} \\
    \hline
    Measurement accuracy & $\pm$0.5\% & $\pm$2.1\% after calibration \\
    \hline
    Bill of materials & Rs.~2,050 & Rs.~1,180 \\
    \hline
    Firmware effort & Low & Moderate \\
    \hline
    Meets cost constraint & No & Yes \\
    \hline
  \end{tabular}
\end{table}

\section{Implementation plan/methodology}

The implementation follows the flow of Figure~\ref{fig:flow}: initialise the
peripherals, acquire one mains cycle of samples, compute RMS quantities,
accumulate energy, decide whether the reading has changed materially, and
publish only when it has.

\begin{figure}[H]
  \centering
  \phbox{3.0in}
  \caption{Implementation flowchart of the firmware}
  \label{fig:flow}
\end{figure}

%------------------------------------------------------------------------------
\chapter{RESULTS ANALYSIS AND VALIDATION}

\section{Implementation of solution}

Modern engineering tools were used throughout:

\begin{itemize}
  \item \textbf{Analysis:} MATLAB was used for sampling-rate and quantisation studies.
  \item \textbf{Design drawings and schematics:} KiCad for schematic capture and PCB layout.
  \item \textbf{Report preparation:} \LaTeX{} for this document; Inkscape for figures.
  \item \textbf{Project management and communication:} Git for version control and a shared Kanban board.
  \item \textbf{Testing and validation:} a calibrated reference meter and a programmable load bank.
\end{itemize}

\subsection{Test setup}

The node was connected in series with a programmable load bank and the reference
meter. Ten load points were applied between 40~W and 2200~W and each point was
held for fifteen minutes. The arrangement is shown in Figure~\ref{fig:setup}.

\begin{figure}[H]
  \centering
  \phbox{2.7in}
  \caption{Experimental test setup used for validation}
  \label{fig:setup}
\end{figure}

\subsection{Measurement results}

Table~\ref{tab:results} lists the measured values against the reference. The
worst-case relative error of 2.1\% occurred at the lowest load point, where the
current transformer operates near the bottom of its linear range.

\begin{table}[H]
  \centering
  \caption{Measured power against the reference meter}
  \label{tab:results}
  \begin{tabular}{|c|c|c|c|}
    \hline
    \textbf{Load point} & \textbf{Reference (W)} & \textbf{Measured (W)} & \textbf{Error (\%)} \\
    \hline
    1 & 40   & 40.8   & 2.10 \\
    \hline
    2 & 250  & 253.1  & 1.24 \\
    \hline
    3 & 800  & 806.4  & 0.80 \\
    \hline
    4 & 1500 & 1508.5 & 0.57 \\
    \hline
    5 & 2200 & 2209.7 & 0.44 \\
    \hline
  \end{tabular}
\end{table}

\subsection{Interpretation and data validation}

Error decreases monotonically with load, which is consistent with the expected
transformer behaviour rather than with a firmware fault. A 72-hour continuous
run produced 51,840 expected records and 51,840 received records, confirming
that the local buffer recovers correctly from short network outages.

%------------------------------------------------------------------------------
\chapter{CONCLUSION AND FUTURE WORK}

\section{Conclusion}

The objective of building a retrofittable, low-cost appliance-level energy
monitoring node was met. The prototype tracked a calibrated reference within
$\pm$2.1\% across the full tested range, against a target of $\pm$3\%, and
dashboard latency stayed under 3~seconds against a target of 5~seconds. The
bill of materials came to Rs.~1,180, comfortably inside the client
cost cap.

The one deviation from the expected result was the error at the lowest load
point, which was higher than the simulation predicted. This is attributed to
core non-linearity in the current transformer below 0.2~A, a region that the
MATLAB model treated as ideal.

\section{Future work}

\begin{itemize}
  \item Replace the current transformer with a Hall-effect sensor to improve low-load linearity.
  \item Extend the firmware to three-phase operation for the hostel distribution boards.
  \item Add a lightweight classifier so a single node can attribute energy to more than one appliance.
  \item Harden the transport with TLS and per-device certificates before wider deployment.
\end{itemize}

%==============================================================================
%  REFERENCES
%==============================================================================
\chapter*{REFERENCES}
\phantomsection\addcontentsline{toc}{chapter}{REFERENCES}

\begin{singlespace}
\begin{enumerate}[leftmargin=1.4cm]
  \item Ariponnammal, S. and Natarajan, S. (1994) `Transport Phenomena of Sm Sel~--~X Asx', Pramana~--~Journal of Physics, Vol.~42, No.~1, pp.~421--425.
  \item Barnard, R.W. and Kellogg, C. (1980) `Applications of Convolution Operators to Problems in Univalent Function Theory', Michigan Mach, J., Vol.~27, pp.~81--94.
  \item Hart, G.W. (1992) `Nonintrusive Appliance Load Monitoring', Proceedings of the IEEE, Vol.~80, No.~12, pp.~1870--1891.
  \item Shin, K.G. and Mckay, N.D. (1984) `Open Loop Minimum Time Control of Mechanical Manipulations and its Applications', Proc.~Amer.~Contr.~Conf., San Diego, CA, pp.~1231--1236.
\end{enumerate}
\end{singlespace}

%==============================================================================
%  APPENDIX
%==============================================================================
\chapter*{APPENDIX}
\phantomsection\addcontentsline{toc}{chapter}{APPENDIX}

\begin{center}\textbf{\fourteen Appendix-1: Bill of Materials}\end{center}

\vspace{6pt}
\begin{singlespace}
\begin{center}
\begin{tabular}{|c|L{6.0cm}|c|c|}
  \hline
  \textbf{S.~No.} & \textbf{Component} & \textbf{Qty} & \textbf{Cost (Rs.)} \\
  \hline
  1 & ESP32 DevKit v1 & 1 & 420 \\
  \hline
  2 & Split-core current transformer, 30~A & 1 & 310 \\
  \hline
  3 & Voltage sensing transformer & 1 & 180 \\
  \hline
  4 & Signal conditioning components & 1 set & 150 \\
  \hline
  5 & Enclosure and hardware & 1 & 120 \\
  \hline
  \multicolumn{3}{|r|}{\textbf{Total}} & \textbf{1180} \\
  \hline
\end{tabular}
\end{center}
\end{singlespace}

%==============================================================================
%  USER MANUAL
%==============================================================================
\chapter*{USER MANUAL}
\phantomsection\addcontentsline{toc}{chapter}{USER MANUAL}

(Complete step by step instructions along with pictures necessary to run the
project.)

\begin{enumerate}[leftmargin=1.4cm]
  \item Switch off the supply at the distribution board before installation.
  \item Clip the split-core current transformer around the live conductor of the
        load to be monitored, observing the arrow direction printed on the core.
  \item Connect the voltage sensing leads to the corresponding terminals of the node.
  \item Restore the supply. The status LED blinks amber while the node searches
        for a network.
  \item Connect a phone to the temporary access point named \texttt{EnergyNode-XXXX}
        and enter the Wi-Fi credentials on the captive portal page.
  \item The LED turns steady green once telemetry is flowing. Readings appear on
        the dashboard within three seconds, as shown in Figure~\ref{fig:dash}.
\end{enumerate}

\begin{figure}[H]
  \centering
  \phbox{2.6in}
  \caption{Dashboard view after successful commissioning}
  \label{fig:dash}
\end{figure}

\end{document}